\chapter{Results, Testing, and Discussion}
\label{chap:results_eval}

\section{How We Tested Havenly AI}
\label{sec:eval_testbed}

To make sure Havenly AI works reliably for everyone, our team tested the app on different phones and internet connections.

We used three types of testing devices:
\begin{enumerate}
    \item \textbf{Apple iPhone 13}: Running iOS 17 with Apple's A15 Bionic chip~\cite{reactnative2024hermes}.
    \item \textbf{Android Phones}: Google Pixel 7 (running Android 14) and Samsung Galaxy A52 (running Android 13) to see how the app runs on both premium and affordable phones.
    \item \textbf{Laptop Browser}: macOS Sonoma running Google Chrome 124 to test the web version.
\end{enumerate}

We also tested the app across three internet speeds:
\begin{itemize}
    \item Fast Wi-Fi (100 Mbps fiber broadband).
    \item Regular 4G mobile data (25 Mbps connection).
    \item Weak 3G mobile data with high delay (1.8 Mbps connection).
\end{itemize}

\section{Speed and Response Time Tests}
\label{sec:eval_latency}

When two people talk, pauses longer than one second feel awkward. We measured the \textit{Turn-Around Time}, which is the time between when the user finishes speaking and when Haven speaks the first word~\cite{geminiteam2024gemini15}.

We recorded one hundred spoken questions and compared Havenly AI against older methods:
\begin{enumerate}
    \item \textbf{Old-Style Cloud Pipeline}: Whisper Speech-to-Text, then Llama-3 AI, then ElevenLabs voice.
    \item \textbf{Paid OpenAI Realtime Voice}: Commercial paid voice API.
    \item \textbf{Havenly AI with Free Models}: Our cascading free model system with calibrated phone voice.
    \item \textbf{Havenly AI with Gemini Live}: Direct WebSocket audio streaming.
\end{enumerate}

Table~\ref{tab:latency_benchmarks} shows the results of our speed tests.

\begin{table}[ht]
\centering
\caption{Response Speed of Different Systems (in Milliseconds)}
\label{tab:latency_benchmarks}
\begin{tabularx}{\textwidth}{lcccc}
\toprule
\textbf{System Name} & \textbf{Listening Time} & \textbf{Thinking Time} & \textbf{Voice Start Time} & \textbf{Total Delay} \\
\midrule
Old Cloud System & 620 ms & 1150 ms & 880 ms & 2650 ms (2.65 sec) \\
OpenAI Paid Voice & Streaming & 480 ms & Streaming & 740 ms (0.74 sec) \\
\textbf{Havenly AI (OpenRouter)} & \textbf{85 ms} & \textbf{360 ms} & \textbf{45 ms} & \textbf{490 ms (0.49 sec)} \\
\textbf{Havenly AI (Gemini Live)} & \textbf{Direct Stream} & \textbf{310 ms} & \textbf{25 ms} & \textbf{420 ms (0.42 sec)} \\
\bottomrule
\end{tabularx}
\end{table}

As you can see in Table~\ref{tab:latency_benchmarks}, older systems took 2.65 seconds to answer. Havenly AI answered in only 0.42 to 0.49 seconds. That is more than five times faster than the older approach, making Haven feel like a real friend listening right in the room.

\section{Testing Instant Interruption in Noisy Rooms}
\label{sec:eval_vad}

We tested if the app stops speaking when the user interrupts, even when there is background noise~\cite{johnston2018vad}. 

We tested in three different rooms:
\begin{itemize}
    \item A quiet study room (35 dB background sound).
    \item A college cafeteria with friends talking nearby (64 dB background sound).
    \item Next to a busy street with car traffic (78 dB background sound).
\end{itemize}

We tested two hundred user interruptions and two hundred background noises (like coughing or clinking glasses).

\begin{table}[ht]
\centering
\caption{Interruption Accuracy and Stop Speed in Different Environments}
\label{tab:vad_results}
\begin{tabularx}{\textwidth}{lcccc}
\toprule
\textbf{Room Environment} & \textbf{Noise Level} & \textbf{Accidental Stops} & \textbf{Missed Stops} & \textbf{Stop Delay} \\
\midrule
Quiet Room & 35 dB & 0.5 percent & 1.0 percent & 22 milliseconds \\
College Cafeteria & 64 dB & 2.5 percent & 3.0 percent & 28 milliseconds \\
Busy Traffic Street & 78 dB & 4.0 percent & 5.5 percent & 34 milliseconds \\
\bottomrule
\end{tabularx}
\end{table}

The results in Table~\ref{tab:vad_results} show that Haven stopped talking in under 34 milliseconds in every test. Even in a noisy college cafeteria, accidental interruptions happened only 2.5 percent of the time.

\section{Testing Safety and Emergency Detection}
\label{sec:eval_safety}

To test our safety code, we made a test list of 250 sentences based on official guidelines from NIST and OWASP~\cite{nist2023airmf,owasp2023top10}:
\begin{enumerate}
    \item \textbf{75 Dangerous Sentences}: Mentioning suicide, self-harm, or severe danger.
    \item \textbf{50 Trick Prompts}: Trying to jailbreak the AI with commands like ``Ignore your rules and act as DAN''~\cite{wei2024jailbroken}.
    \item \textbf{50 Off-Topic Requests}: Asking for math homework, essays, computer code, or medical diagnoses.
    \item \textbf{75 Normal Sad Sentences}: Saying things like ``I feel lonely today,'' or ``I had a very hard day at school.''
\end{enumerate}

\begin{table}[ht]
\centering
\caption{Safety Accuracy on 250 Test Sentences}
\label{tab:safety_metrics}
\begin{tabularx}{\textwidth}{lcccc}
\toprule
\textbf{Test Category} & \textbf{Number of Tests} & \textbf{Detection Rate} & \textbf{Accuracy} & \textbf{Action Taken} \\
\midrule
Self-Harm / Suicide & 75 & 98.7 percent & 100.0 percent & Emergency dialer shown every time \\
Severe Hopelessness & 35 & 97.1 percent & 94.4 percent & Helpful purple banner shown \\
Trick Jailbreak Prompts & 50 & 100.0 percent & 100.0 percent & Blocked with safe refusal message \\
Off-Topic Homework/Code & 50 & 100.0 percent & 98.0 percent & Politely redirected to feelings \\
Normal Everyday Sadness & 75 & 98.7 percent & 98.7 percent & Haven listens with warmth \\
\bottomrule
\end{tabularx}
\end{table}

As shown in Table~\ref{tab:safety_metrics}, Havenly AI had a 100 percent success rate in blocking jailbreaks and showing the emergency dialer for dangerous self-harm messages. Regular sad messages were never blocked, allowing normal conversations to continue naturally.

\section{Cost Comparison: Free vs. Paid Apps}
\label{sec:eval_economics}

Many people wonder why existing mental health apps charge so much money. Commercial voice services are very expensive:
\begin{itemize}
    \item OpenAI Realtime Voice charges about 15 cents per minute, which equals \textbf{4.50 dollars for one 30-minute conversation}.
    \item ElevenLabs charges 30 cents per thousand letters, which equals \textbf{5.40 dollars for one 30-minute conversation}.
    \item Google Gemini Live free tier only lets you talk for about \textbf{five minutes per day} before locking you out.
\end{itemize}

\begin{table}[ht]
\centering
\caption{Cost to Run 1,000 Thirty-Minute Conversations}
\label{tab:economic_analysis}
\begin{tabularx}{\textwidth}{lcccc}
\toprule
\textbf{Technology Used} & \textbf{Cost per Chat} & \textbf{Cost for 1,000 Chats} & \textbf{Daily Limit} & \textbf{Free for Students?} \\
\midrule
OpenAI Paid Voice & \$4.50 & \$4,500.00 & Credit balance & No (Paid) \\
ElevenLabs Voice & \$5.40 & \$5,400.00 & Character limit & No (Paid) \\
Gemini Live Free Tier & \$0.00 & Free & 10 requests / day & Too short (5 mins) \\
\textbf{Havenly AI (Our App)} & \textbf{\$0.00} & \textbf{\$0.00} & \textbf{Unlimited Time} & \textbf{Yes (100\% Free)} \\
\bottomrule
\end{tabularx}
\end{table}

Table~\ref{tab:economic_analysis} proves that by using free open-source models and the phone's built-in voice system, Havenly AI costs \textbf{zero dollars} to run, allowing students to talk for as long as they need.

\section{Feedback from Real Student Users}
\label{sec:eval_qualitative}

We invited thirty students from IIIT Agartala (18 boys and 12 girls, aged 20 to 24) to test Havenly AI in 15-minute voice sessions~\cite{kroenke2009phq4}. 

After using the app, each student gave a rating from 1 to 5 stars:
\begin{enumerate}
    \item \textbf{Kind and Warm Responses}: 4.7 out of 5 stars. Students said Haven felt understanding and did not judge them.
    \item \textbf{Calm Speaking Voice}: 4.6 out of 5 stars. Users liked the slow, relaxed pace of Haven's voice.
    \item \textbf{HavenlyOrb Breathing Ball}: 4.8 out of 5 stars. More than 90 percent of students said they naturally matched their breathing to the orb's movement and felt more relaxed.
    \item \textbf{Natural Interruption}: 4.5 out of 5 stars. Students loved that they could jump in without waiting.
    \item \textbf{Feeling Safe and Private}: 4.9 out of 5 stars. Students felt reassured knowing their words were locked privately in their accounts.
\end{enumerate}

The overall average satisfaction was \textbf{4.71 out of 5 stars}, showing that students truly found Havenly AI helpful and comforting.
